% Lecture 4 -- Ito formula: intuition, rules, and applications
\documentclass[10pt,aspectratio=169]{beamer}

\usepackage[utf8]{inputenc}
\usepackage[T1]{fontenc}
\usepackage[english]{babel}
\usepackage{lmodern}
\usepackage{microtype}
\usepackage{amsmath,amssymb,bm,mathtools}
\usepackage{booktabs}
\usepackage{tikz}
\usetikzlibrary{arrows.meta,positioning,calc}

\usetheme[numbering=fraction,progressbar=frametitle]{metropolis}
\setbeamertemplate{frame numbering}[fraction]
\setbeamertemplate{navigation symbols}{}

\definecolor{BootcampBlueDark}{HTML}{1A3C68}
\definecolor{BootcampBlue}{RGB}{31,110,178}
\definecolor{AccentGold}{HTML}{DAA520}
\definecolor{SoftGray}{RGB}{245,246,248}
\definecolor{TextBlack}{HTML}{000000}
\definecolor{Crimson}{HTML}{990000}
\definecolor{MatrixGreen}{HTML}{228B22}

\setbeamercolor{normal text}{fg=TextBlack,bg=white}
\setbeamercolor{title}{fg=TextBlack}
\setbeamerfont{title}{size=\Large,series=\bfseries}
\setbeamercolor{structure}{fg=BootcampBlueDark}
\setbeamercolor{progress bar}{fg=BootcampBlueDark,bg=gray!30}
\setbeamercolor{frametitle}{bg=BootcampBlueDark,fg=white}
\setbeamerfont{frametitle}{series=\bfseries,size=\large}
\setbeamertemplate{frametitle}{
  \nointerlineskip
  \begin{beamercolorbox}[wd=\paperwidth,ht=3.2ex,dp=1.4ex,leftskip=1em,rightskip=1em]{frametitle}
    \usebeamerfont{frametitle}\insertframetitle
  \end{beamercolorbox}
}

\setbeamercolor{block title example}{bg=BootcampBlueDark!90!black,fg=white}
\setbeamercolor{block body example}{bg=BootcampBlueDark!5!white,fg=black}
\setbeamercolor{block title proposition}{bg=MatrixGreen!75!black,fg=white}
\setbeamercolor{block body proposition}{bg=MatrixGreen!7!white,fg=black}
\setbeamertemplate{blocks}[rounded][shadow=false]

\newenvironment{definitionblock}[1]{%
  \par\vspace{0.35em}%
  \begin{beamercolorbox}[wd=\textwidth,sep=0.35em,leftskip=0.3em,rightskip=0.6em]{block title example}%
    \raisebox{0.12em}{\usebeamerfont{block title example}#1}%
  \end{beamercolorbox}%
  \vspace{-0.35em}%
  \begin{beamercolorbox}[wd=\textwidth,sep=0.75em,leftskip=0.6em,rightskip=0.6em]{block body example}%
    \usebeamerfont{block body example}%
}{%
  \end{beamercolorbox}%
  \par\vspace{0.45em}%
}

\newenvironment{propositionblock}[1]{%
  \par\vspace{0.35em}%
  \begin{beamercolorbox}[wd=\textwidth,sep=0.35em,leftskip=0.3em,rightskip=0.6em]{block title proposition}%
    \raisebox{0.12em}{\usebeamerfont{block title example}#1}%
  \end{beamercolorbox}%
  \vspace{-0.35em}%
  \begin{beamercolorbox}[wd=\textwidth,sep=0.75em,leftskip=0.6em,rightskip=0.6em]{block body proposition}%
    \usebeamerfont{block body example}%
}{%
  \end{beamercolorbox}%
  \par\vspace{0.45em}%
}

\newcommand{\E}{\mathbb E}
\newcommand{\Pp}{\mathbb P}
\newcommand{\F}{\mathcal F}
\newcommand{\R}{\mathbb R}
\newcommand{\dd}{\,\mathrm d}
\newcommand{\Ito}{\mathrm{It\hat{o}}}
\DeclareMathOperator{\Var}{Var}
\DeclareMathOperator{\Cov}{Cov}

% One clean transition page per major section.
\AtBeginSection[]{
  \begin{frame}[plain]
    \vfill
    {\Large\bfseries\color{BootcampBlueDark}\insertsectionhead\par}
    \medskip
    {\color{BootcampBlueDark}\rule{0.9\textwidth}{0.8pt}}
    \vfill
  \end{frame}
}

\title{The It\^o Formula: Intuition, Rules, and Applications}
\author{}
\date{}

\begin{document}

\begin{frame}[plain]
  \vfill
  {\usebeamerfont{title}\usebeamercolor[fg]{title}\inserttitle\par}
  \medskip
  {\color{BootcampBlueDark}\rule{0.95\textwidth}{0.8pt}}
  \vfill
\end{frame}

\begin{frame}{References}
\footnotesize
\textbf{Main references}
\begin{itemize}\setlength{\itemsep}{0.55em}
  \item S.~Shreve, \emph{Stochastic Calculus for Finance II}, Chapter 4.
  \item I.~Karatzas and S.~Shreve, \emph{Brownian Motion and Stochastic Calculus}, chapter on stochastic integration and \(\Ito\)'s rule.
  \item T.~Bj\"ork, \emph{Arbitrage Theory in Continuous Time}, chapters on stochastic calculus and stochastic differential equations.
\end{itemize}
\begin{block}{Scope of this lecture}
We learn the second-order chain rule, use it to solve basic SDEs, and connect it to products, correlation, discounting, and hedging.
\end{block}
\end{frame}

\begin{frame}{Roadmap}
\begin{enumerate}\setlength{\itemsep}{0.48em}
  \item See why the ordinary chain rule breaks down for a random path, then
        turn quadratic variation into a small multiplication table.
  \item Understand and apply the \(\Ito\) formula for Brownian motion.
  \item Extend the rule to general one-dimensional \(\Ito\) processes.
  \item Solve GBM and OU/Vasicek, then obtain the mean of a CIR process.
  \item Handle products, ratios, and correlated risk factors.
  \item Finish with discounting, convexity, and delta-hedging intuition.
\end{enumerate}
\end{frame}

% =====================================================
\section{Quadratic Variation Becomes a Calculus}
\subsection{Motivation: why the chain rule needs fixing}

\begin{frame}{Why do we even need a new chain rule?}
In finance we constantly want to know how a value that depends on the
stock, \(V_t=v(t,S_t)\), changes as \(S_t\) moves.
\begin{block}{The tool we already know}
Ordinary calculus says
\[
\dd v(t,S_t)=v_t\,\dd t+v_s\,\dd S_t.
\]
\end{block}
\begin{alertblock}{The question for today}
Stock prices do not move along smooth, differentiable paths. Does this
formula still hold when \(S_t\) is driven by Brownian motion?
\end{alertblock}
\end{frame}

\begin{frame}{Let's just try it, and see what happens}
Take the simplest possible case, \(f(x)=x^2\) applied to Brownian motion
itself. If the ordinary chain rule were correct, we would guess
\[
\dd(W_t^2)\overset{?}{=}2W_t\dd W_t,
\qquad\text{i.e.}\qquad
W_T^2\overset{?}{=}2\int_0^T W_t\dd W_t.
\]
\begin{block}{Your turn}
The right-hand side is a stochastic integral with mean zero. Take
expectations on both sides. Does the result agree with what you already
know about \(\E[W_T^2]\)?
\end{block}
\end{frame}

\begin{frame}{The naive guess is simply wrong}
Taking expectations of the naive guess gives
\[
\E[W_T^2]\overset{?}{=}2\,\E\!\left[\int_0^TW_t\dd W_t\right]=0.
\]
But we already know
\[
\boxed{\E[W_T^2]=T\neq0.}
\]
\begin{alertblock}{Something is missing}
The ordinary chain rule cannot be correct for Brownian motion. An entire
term of order \(T\) has been dropped.
\end{alertblock}
\end{frame}

\begin{frame}{Where a smooth path would give zero, Brownian motion does not}
For a smooth, differentiable path \(x(t)\), a small move satisfies
\(\Delta x\approx x'(t)\Delta t\), so \((\Delta x)^2\) is of order
\((\Delta t)^2\): negligible after summation.
\begin{block}{Brownian motion is rougher}
Brownian increments scale like
\[
\Delta W\sim\sqrt{\Delta t}
\quad\Longrightarrow\quad
(\Delta W)^2\sim\Delta t.
\]
\end{block}
Squared increments of a smooth path vanish in the limit; squared increments
of Brownian motion do not. That is the entire source of the missing term.
\end{frame}

\begin{frame}{A quick numerical check}
Split \([0,T]\) into \(n\) equal steps and add up the squared increments
\(\sum_i(\Delta W_i)^2\). Simulating this on a computer for growing \(n\)
gives, for \(T=1\):
\begin{center}
\begin{tabular}{cccc}
\toprule
\(n\) & 10 & 1{,}000 & 100{,}000 \\
\midrule
\(\sum_i(\Delta W_i)^2\) & \(\approx0.8\) & \(\approx0.97\) & \(\approx1.00\) \\
\bottomrule
\end{tabular}
\end{center}
\begin{block}{The pattern}
Regardless of the random path drawn, the sum keeps converging to \(T=1\).
This limit is exactly the quadratic variation \([W]_T\) from Lecture~2.
\end{block}
\end{frame}

\begin{frame}{This is why we revisit quadratic variation today}
We are not repeating Lecture~2 for its own sake. We now know precisely why
it matters:
\begin{enumerate}
  \item the ordinary chain rule silently assumes squared increments vanish;
  \item for Brownian motion they do not --- they accumulate to \([W]_T=T\);
  \item so the chain rule needs a correction term built from quadratic
        variation.
\end{enumerate}
\begin{block}{Today's goal}
Turn ``squared increments accumulate to \(\dd t\)'' into a systematic
bookkeeping rule, then use it to fix the chain rule once and for all.
\end{block}
\end{frame}

\subsection{From scales to multiplication rules}

\begin{frame}{One identity from last lecture contains the new calculus}
Left-point sums and Brownian quadratic variation gave
\[
\int_0^T W_t\dd W_t=\frac12(W_T^2-T).
\]
Rearranging,
\[
\boxed{W_T^2=2\int_0^T W_t\dd W_t+T.}
\]
\begin{columns}[T,onlytextwidth]
\column{0.48\textwidth}
\begin{block}{Ordinary-looking part}
The term \(2W_t\dd W_t\) is what the classical chain rule predicts.
\end{block}
\column{0.48\textwidth}
\begin{block}{New part}
The extra \(T\) comes from \([W]_T=T\).
\end{block}
\end{columns}
Today we turn that correction into a systematic rule.
\end{frame}

\begin{frame}{The classical chain rule misses an order-one term}
On one small interval,
\[
(W_t+\Delta W)^2-W_t^2
=2W_t\Delta W+(\Delta W)^2.
\]
For a smooth path, the square of the increment disappears after summation. For Brownian motion,
\[
\Delta W\sim\sqrt{\Delta t}
\quad\Longrightarrow\quad
(\Delta W)^2\sim\Delta t.
\]
Thus \(n\) small quadratic terms accumulate to a finite correction.
\begin{alertblock}{The key change}
First-order calculus is not enough: a second-order Taylor term survives.
\end{alertblock}
\end{frame}

\begin{frame}{Orders of magnitude tell us which products survive}
On a uniform grid with \(n\) steps, \(\Delta t=T/n\):
\begin{center}
\begin{tabular}{lccc}
\toprule
One-step term & Typical size & Sum of \(n\) terms & Limit scale \\
\midrule
\(\Delta t\) & \(n^{-1}\) & \(1\) & survives \\
\((\Delta W)^2\) & \(n^{-1}\) & \(1\) & survives \\
\(\Delta t\,\Delta W\) & \(n^{-3/2}\) & \(n^{-1/2}\) & vanishes \\
\((\Delta t)^2\) & \(n^{-2}\) & \(n^{-1}\) & vanishes \\
\bottomrule
\end{tabular}
\end{center}
\begin{block}{Bookkeeping principle}
When expanding stochastic differentials, retain terms that accumulate at order \(\dd t\).
\end{block}
\end{frame}

\begin{frame}{Covariation records accumulated products of increments}
For continuous processes \(X,Y\) and a partition \(\pi\), set
\[
[X,Y]_t^\pi
=\sum_i(X_{t_{i+1}}-X_{t_i})(Y_{t_{i+1}}-Y_{t_i}).
\]
When the limit exists as \(|\pi|\to0\), it is denoted by \([X,Y]_t\).
\begin{propositionblock}{Basic examples}
\[
[W,W]_t=t,
\qquad
[W,t]_t=0,
\qquad
[t,t]_t=0.
\]
\end{propositionblock}
Quadratic variation is the special case \([X]_t=[X,X]_t\).
\end{frame}

\begin{frame}{The one-factor multiplication table}
The differential notation summarizes the covariation limits:
\begin{center}
\renewcommand{\arraystretch}{1.55}
\begin{tabular}{c|cc}
\toprule
\(\times\) & \(\dd t\) & \(\dd W_t\) \\
\midrule
\(\dd t\) & \(0\) & \(0\) \\
\(\dd W_t\) & \(0\) & \(\dd t\) \\
\bottomrule
\end{tabular}
\end{center}
Equivalently,
\[
(\dd t)^2=0,
\qquad
\dd t\,\dd W_t=0,
\qquad
(\dd W_t)^2=\dd t.
\]
\begin{alertblock}{How to read the symbols}
The product \(\dd X_t\,\dd Y_t\) is shorthand for the covariation increment
\(\dd[X,Y]_t\). The table is a mnemonic for limits, not literal infinitesimal algebra.
\end{alertblock}
\end{frame}

\begin{frame}{The table with two Brownian shocks includes correlation}
Let \(W\) and \(B\) be Brownian motions with
\[
[W,B]_t=\rho t,
\qquad -1\le\rho\le1.
\]
Then
\begin{center}
\renewcommand{\arraystretch}{1.5}
\begin{tabular}{c|ccc}
\toprule
\(\times\) & \(\dd t\) & \(\dd W_t\) & \(\dd B_t\) \\
\midrule
\(\dd t\) & \(0\) & \(0\) & \(0\) \\
\(\dd W_t\) & \(0\) & \(\dd t\) & \(\rho\,\dd t\) \\
\(\dd B_t\) & \(0\) & \(\rho\,\dd t\) & \(\dd t\) \\
\bottomrule
\end{tabular}
\end{center}
For independent Brownian motions, \(\rho=0\) and \(\dd W_t\dd B_t=0\).
More generally, \(\dd W_t^i\dd W_t^j=\rho_{ij}\dd t\) for a Brownian vector with correlation matrix \((\rho_{ij})\).
\end{frame}

\begin{frame}{Application: correlation appears as drift in a product}
Apply the discrete product identity in differential form:
\[
\dd(W_tB_t)=W_t\dd B_t+B_t\dd W_t+\dd W_t\dd B_t.
\]
Using the table,
\[
\boxed{\dd(W_tB_t)=W_t\dd B_t+B_t\dd W_t+\rho\dd t.}
\]
After integration and expectation,
\[
\E[W_tB_t]=\rho t.
\]
\begin{block}{Sanity check}
The stochastic integrals have mean zero under their \(L^2\) conditions, leaving exactly the covariance \(\rho t\).
\end{block}
\end{frame}

\begin{frame}{Table drill: same shock, independent shock, correlated shock}
Let
\[
\dd X_t=2\dd t+3\dd W_t,
\quad
\dd Y_t=-\dd t+4\dd W_t,
\quad
\dd Z_t=-\dd t+4\dd B_t.
\]
\begin{block}{Your turn}
\begin{enumerate}\setlength{\itemsep}{0.5em}
  \item Compute \((\dd X_t)^2\).
  \item Compute \(\dd X_t\dd Y_t\): both processes use the same shock.
  \item Compute \(\dd X_t\dd Z_t\) when \(W,B\) are independent.
  \item Repeat when \(\dd W_t\dd B_t=-\tfrac12\dd t\).
\end{enumerate}
\end{block}
Predict the sign before doing the last calculation.
\end{frame}

\begin{frame}{Table drill debrief: only diffusion coefficients survive}
All terms containing \(\dd t\) disappear from products of differentials:
\[
\boxed{
(\dd X_t)^2=9\dd t,
\qquad
\dd X_t\dd Y_t=12\dd t,
\qquad
\dd X_t\dd Z_t=12\rho\dd t.}
\]
Therefore
\[
\dd X_t\dd Z_t=
\begin{cases}
0, & \rho=0,\\
-6\dd t, & \rho=-\tfrac12.
\end{cases}
\]
\begin{block}{Takeaway}
The coefficients \(2,-1\) multiplying \(\dd t\) never enter covariation. Only diffusion loadings and correlation matter.
\end{block}
\end{frame}

% =====================================================
\section{The Ito Formula Is a Second-Order Chain Rule}
\subsection{Statement, construction, and first applications}

\begin{frame}{Brownian scaling tells us how far Taylor must go}
For a smooth function \(f(t,x)\), write heuristically
\[
\Delta f
\approx f_t\Delta t+f_x\Delta W
+\frac12f_{xx}(\Delta W)^2.
\]
The omitted terms have smaller accumulated order:
\[
\Delta t\,\Delta W\sim(\Delta t)^{3/2},
\qquad
(\Delta W)^3\sim(\Delta t)^{3/2}.
\]
Replacing \((\Delta W)^2\) by \(\Delta t\) suggests
\[
\Delta f\approx
\left(f_t+\frac12f_{xx}\right)\Delta t+f_x\Delta W.
\]
\end{frame}

\begin{frame}{The \(\Ito\) formula for \(f(t,W_t)\)}
Let \(f\in C^{1,2}([0,T]\times\R)\).
\begin{propositionblock}{Brownian \(\Ito\) formula}
\begin{align*}
f(t,W_t)=f(0,W_0)
&+\int_0^t f_t(s,W_s)\dd s
+\int_0^t f_x(s,W_s)\dd W_s\\
&+\frac12\int_0^t f_{xx}(s,W_s)\dd s.
\end{align*}
\end{propositionblock}
In differential notation,
\[
\boxed{\dd f(t,W_t)
=\left(f_t+\tfrac12f_{xx}\right)(t,W_t)\dd t
+f_x(t,W_t)\dd W_t.}
\]
The term involving $dt$ is called the drift, while the term involving $dW_t$ is called the diffusion (or volatility) term.
\end{frame}

\begin{frame}{Construction: telescope, expand, and keep surviving scales}
On a partition \(0=t_0<\cdots<t_n=t\), telescope
\[
f(t,W_t)-f(0,W_0)=\sum_i\Delta f_i.
\]
A second-order Taylor expansion gives
\[
\Delta f_i
=f_t\Delta t_i+f_x\Delta W_i
+\frac12f_{xx}(\Delta W_i)^2+r_i.
\]
As the mesh vanishes:
\begin{itemize}\setlength{\itemsep}{0.4em}
  \item the first sum becomes \(\int f_t\dd t\);
  \item the left-point second sum becomes \(\int f_x\dd W\);
  \item quadratic variation turns the third sum into \(\frac12\int f_{xx}\dd t\);
  \item the accumulated remainder vanishes after localization.
\end{itemize}
\end{frame}

\begin{frame}{A four-step recipe prevents most mistakes}
To compute \(\dd f(t,W_t)\):
\begin{enumerate}\setlength{\itemsep}{0.65em}
  \item Treat \(t\) and \(x\) as ordinary deterministic variables.
  \item Compute \(f_t(t,x)\), \(f_x(t,x)\), and \(f_{xx}(t,x)\).
  \item Evaluate the derivatives at \(x=W_t\).
  \item Assemble the \(\dd t\) and \(\dd W_t\) terms separately.
\end{enumerate}
\begin{alertblock}{Notation check}
Differentiate with respect to \(x\), not ``with respect to \(W_t\).'' Compute first, then substitute \(x=W_t\).
\end{alertblock}
\end{frame}

\begin{frame}{Worked examples: powers and a curved transform}
For \(f(x)=x^2\),
\[
\boxed{\dd(W_t^2)=2W_t\dd W_t+\dd t.}
\]
Moreover, on \([0,T]\),
\(
\E\int_0^T4W_s^2\dd s=2T^2<\infty
\), so \(W_t^2-t\) is a true martingale.

For \(f(x)=x^3\),
\[
\boxed{\dd(W_t^3)=3W_t^2\dd W_t+3W_t\dd t.}
\]
For \(f(x)=\sin x\),
\[
\boxed{\dd(\sin W_t)=\cos(W_t)\dd W_t-\tfrac12\sin(W_t)\dd t.}
\]
\begin{block}{What changed?}
The curvature \(f''\) creates drift even though Brownian motion itself has zero drift.
\end{block}
\end{frame}

\begin{frame}{Exercise: manufacture a fourth-order martingale}
Apply \(\Ito\)'s formula to \(f(x)=x^4\).
\begin{block}{Your turn}
\begin{enumerate}\setlength{\itemsep}{0.55em}
  \item Compute \(\dd(W_t^4)\).
  \item Find \(c\) such that
  \[
  W_t^4-c\int_0^tW_s^2\dd s
  \]
  has zero drift.
  \item Use expectations to recover \(\E[W_t^4]\).
\end{enumerate}
\end{block}
Why is the resulting local martingale a true martingale on every bounded horizon?
\end{frame}

\begin{frame}{Exercise debrief: curvature generates the moment recursion}
Since \(f'(x)=4x^3\) and \(f''(x)=12x^2\),
\[
\dd(W_t^4)=4W_t^3\dd W_t+6W_t^2\dd t.
\]
Therefore
\[
\boxed{M_t=W_t^4-6\int_0^tW_s^2\dd s}
\]
is a true martingale on bounded horizons; the polynomial integrands have all required moments. Taking expectations,
\[
\E[W_t^4]
=6\int_0^t\E[W_s^2]\dd s
=6\int_0^ts\dd s
=3t^2.
\]
\end{frame}

\begin{frame}{Exercise: one more step gives the sixth moment}
Fix \(T>0\) and use the result \(\E[W_t^4]=3t^2\).
\begin{block}{Your turn}
\begin{enumerate}\setlength{\itemsep}{0.55em}
  \item Apply \(\Ito\)'s formula to \(f(x)=x^6\).
  \item Separate the drift and stochastic-integral terms.
  \item Take expectations and compute \(\E[W_T^6]\).
\end{enumerate}
\end{block}
Why is the stochastic integral a true martingale on every bounded horizon?
\end{frame}

\begin{frame}{Exercise debrief: the moment recursion closes}
Since \(f'(x)=6x^5\) and \(f''(x)=30x^4\),
\[
\dd(W_t^6)=6W_t^5\dd W_t+15W_t^4\dd t.
\]
Because \(\E\int_0^T36W_t^{10}\dd t<\infty\), the stochastic integral is a true martingale on \([0,T]\). Hence
\[
\boxed{
\E[W_T^6]
=15\int_0^T\E[W_t^4]\dd t
=45\int_0^Tt^2\dd t
=15T^3.}
\]
More generally, \(m_n(t)=\E[W_t^n]\) satisfies
\[
m_n'(t)=\frac{n(n-1)}2m_{n-2}(t).
\]
\end{frame}

\begin{frame}{Exercise: exit time and occupation inside an interval}
For \(a>0\), let
\[
\tau_a=\inf\{t\ge0:|W_t|=a\}.
\]
Use the martingales
\[
W_t^2-t,
\qquad
W_t^4-6\int_0^tW_s^2\dd s.
\]
\begin{block}{Your turn}
Stop first at \(t\wedge\tau_a\), then determine
\[
\E[\tau_a]
\qquad\text{and}\qquad
\E\!\left[\int_0^{\tau_a}W_s^2\dd s\right].
\]
\end{block}
Why should the limit \(t\to\infty\) be taken only after stopping?
\end{frame}

\begin{frame}{Exercise debrief: stopping turns polynomials into identities}
Write \(\tau_t=t\wedge\tau_a\). Optional sampling gives
\[
\E[\tau_t]=\E[W_{\tau_t}^2]\le a^2.
\]
Monotone convergence first yields \(\E[\tau_a]\le a^2\), hence \(\tau_a<\infty\) a.s. Then continuity and dominated convergence give
\[
\boxed{\E[\tau_a]=a^2.}
\]
Stopping the fourth-order martingale similarly gives
\[
\E[W_{\tau_t}^4]
=6\E\!\left[\int_0^{\tau_t}W_s^2\dd s\right].
\]
Letting \(t\to\infty\) now yields
\[
\boxed{\E\!\left[\int_0^{\tau_a}W_s^2\dd s\right]=\frac{a^4}{6}.}
\]
\end{frame}

\begin{frame}{The corrected exponential has exactly zero drift}
Let
\[
M_t=\exp\left(\lambda W_t-\frac12\lambda^2t\right).
\]
For \(f(t,x)=e^{\lambda x-\lambda^2t/2}\),
\[
f_t=-\frac12\lambda^2f,
\qquad f_x=\lambda f,
\qquad f_{xx}=\lambda^2f.
\]
Thus the two drift terms cancel:
\[
\boxed{\dd M_t=\lambda M_t\dd W_t.}
\]
Here \(M_t\ge0\) and direct Gaussian calculation gives \(\E[M_t]=1\), so \(M\) is a true martingale.
\end{frame}

\begin{frame}{Zero drift is a differential equation for the function}
For \(Y_t=u(t,W_t)\), the \(\Ito\) formula gives
\[
\dd Y_t
=\underbrace{\left(u_t+\frac12u_{xx}\right)(t,W_t)\dd t}_{\text{drift}}
+u_x(t,W_t)\dd W_t.
\]
Thus the drift vanishes when
\[
\boxed{u_t+\frac12u_{xx}=0.}
\]
\begin{block}{Probabilistic reading}
Solutions of the backward heat equation turn Brownian motion into local martingales. Suitable growth or integrability conditions upgrade them to true martingales.
\end{block}
\end{frame}

\begin{frame}{Exercise: which transforms have zero drift?}
For fixed \(T>0\) and \(\lambda\in\R\), consider
\[
u_1(t,x)=x^2-t,
\quad
u_2(t,x)=x^2+t,
\quad
u_3(t,x)=x^2+T-t,
\quad
u_4(t,x)=e^{\lambda x-\lambda^2t/2}.
\]
\begin{block}{Your turn}
For each function, compute \(u_t+\frac12u_{xx}\). Decide whether \(u_i(t,W_t)\) has zero drift and interpret \(u_3(t,W_t)\).
\end{block}
\end{frame}

\begin{frame}{Exercise debrief: the generator gives an instant test}
\begin{center}
\begin{tabular}{lcc}
\toprule
Function & \(u_t+\frac12u_{xx}\) & Conclusion \\
\midrule
\(x^2-t\) & \(0\) & martingale \\
\(x^2+t\) & \(2\) & drift \(2\dd t\) \\
\(x^2+T-t\) & \(0\) & martingale \\
\(e^{\lambda x-\lambda^2t/2}\) & \(0\) & martingale \\
\bottomrule
\end{tabular}
\end{center}
Moreover,
\[
W_t^2+T-t=\E[W_T^2\mid\F_t].
\]
\begin{block}{Takeaway}
The backward equation encodes a conditional forecast that stays unchanged on average as information arrives.
\end{block}
\end{frame}

% ================================================================
\section{From Brownian Motion to Ito Processes}
\subsection{General formula, moments, and logarithms}
% ================================================================

\begin{frame}{An Ito process separates trend from noise}
On a fixed horizon \([0,T]\), an \emph{Ito process} has the form
\[
X_t=X_0+\int_0^t a_s\dd s+\int_0^t b_s\dd W_s.
\]
We assume, for instance,
\[
\int_0^T|a_s|\dd s<\infty\quad\text{a.s.},
\qquad
\E\!\left[\int_0^T b_s^2\dd s\right]<\infty,
\]
with progressively measurable coefficients \(a\) and \(b\).
\begin{center}
\begin{tabular}{ccc}
\toprule
Term & Role & Size over \(\dd t\) \\
\midrule
\(a_t\dd t\) & local trend & \(\dd t\) \\
\(b_t\dd W_t\) & local shock & \(\sqrt{\dd t}\) \\
\bottomrule
\end{tabular}
\end{center}
The shorthand is \(\dd X_t=a_t\dd t+b_t\dd W_t\).
\end{frame}

\begin{frame}{Differential and integral forms describe the same process}
\small
The differential form is compact local notation:
\[
\boxed{\dd X_t=a_t\dd t+b_t\dd W_t.}
\]
Its precise integral meaning, on every interval \(0\leq s\leq t\leq T\), is
\[
\boxed{
X_t-X_s=\int_s^t a_u\dd u+\int_s^t b_u\dd W_u.}
\]
\begin{columns}[T,onlytextwidth]
\column{0.47\textwidth}
\begin{block}{Differential form}
Best for local calculations, It\^o's formula, and identifying drift and
diffusion.
\end{block}
\column{0.49\textwidth}
\begin{exampleblock}{Integral form}
Best for stating what the process actually is and comparing values at two
times.
\end{exampleblock}
\end{columns}
\begin{alertblock}{Do not divide by \(\dd t\)}
Brownian paths are not differentiable: \(\dd W_t/\dd t\) is not an ordinary
random function. The differential equation must be understood through its
integral form.
\end{alertblock}
\end{frame}

\begin{frame}{Optional, for the curious: the drift and diffusion you read off are the only ones}
\small
If you write the same process two different ways,
\[
X_t=X_0+\int_0^t a_s\dd s+\int_0^t b_s\dd W_s
    =X_0+\int_0^t \widetilde a_s\dd s
          +\int_0^t \widetilde b_s\dd W_s,
\]
it turns out \(a=\widetilde a\) and \(b=\widetilde b\) --- there is only one
correct drift and only one correct diffusion coefficient for a given
process. Intuitively: a purely-\(\dd t\) term (finite variation, zero
quadratic variation) and a purely-\(\dd W\) term (zero drift on average,
positive quadratic variation) have such different ``signatures'' that they
can only cancel each other if both are zero.
\begin{block}{Why this matters in practice}
It means the drift/diffusion split you get after applying \(\Ito\)'s
formula is not a matter of algebraic convenience: it is \emph{the} answer,
and every later formula (delta hedge, PDE, martingale test) can safely rely
on it.
\end{block}
\end{frame}

\begin{frame}{Quadratic variation sees volatility, not drift}
If
\[
\dd X_t=a_t\dd t+b_t\dd W_t,
\]
then the multiplication table gives
\[
(\dd X_t)^2
=a_t^2(\dd t)^2+2a_tb_t\dd t\dd W_t+b_t^2(\dd W_t)^2
=b_t^2\dd t.
\]
Therefore
\[
\boxed{\dd[X]_t=b_t^2\dd t,\qquad [X]_t=\int_0^t b_s^2\dd s.}
\]
\begin{block}{Interpretation}
Quadratic variation measures accumulated instantaneous variance. A finite-variation drift contributes nothing to it.
\end{block}
\end{frame}

\begin{frame}{The general one-dimensional Ito formula}
Let
\[
\dd X_t=a_t\dd t+b_t\dd W_t,
\qquad f\in C^{1,2}.
\]
Then
\[
\boxed{
\dd f(t,X_t)
=\left(f_t+a_tf_x+\frac12b_t^2f_{xx}\right)(t,X_t)\dd t
+b_tf_x(t,X_t)\dd W_t.}
\]
Equivalently,
\[
\dd f=f_t\dd t+f_x\dd X+\frac12f_{xx}\dd[X].
\]
\begin{block}{What changed from ordinary calculus?}
Only the curvature correction \(\frac12f_{xx}\dd[X]\). Everything else looks like the usual chain rule.
\end{block}
\end{frame}

\begin{frame}{Worked example: square a general Ito process}
Apply \(f(x)=x^2\) to \(\dd X_t=a_t\dd t+b_t\dd W_t\):
\[
f'(x)=2x,\qquad f''(x)=2.
\]
Hence
\[
\boxed{\dd(X_t^2)=(2X_ta_t+b_t^2)\dd t+2X_tb_t\dd W_t.}
\]
The term \(b_t^2\dd t\) is invisible to the classical chain rule.
\begin{propositionblock}{Compensated square}
Whenever the stochastic integral is integrable enough,
\[
X_t^2-X_0^2-\int_0^t(2X_sa_s+b_s^2)\dd s
\]
is a martingale.
\end{propositionblock}
\end{frame}

\begin{frame}{Exercise: arithmetic diffusion and its second moment}
Let
\[
X_t=x+\mu t+\sigma W_t.
\]
\begin{block}{Your turn}
\begin{enumerate}
\item Compute \(\dd(X_t^2)\).
\item Find a compensated process with zero drift.
\item Use the result to calculate \(\E[X_t^2]\).
\end{enumerate}
\end{block}
\begin{alertblock}{Integrability check}
Before taking expectations, explain why the stochastic integral has mean zero on \([0,T]\).
\end{alertblock}
\end{frame}

\begin{frame}{Exercise debrief: noise adds variance to the square}
Since \(\dd X_t=\mu\dd t+\sigma\dd W_t\),
\[
\dd(X_t^2)=(2\mu X_t+\sigma^2)\dd t+2\sigma X_t\dd W_t.
\]
Thus
\[
M_t=X_t^2-x^2-\int_0^t(2\mu X_s+\sigma^2)\dd s
\]
has zero drift. On bounded horizons, \(\E\int_0^T X_s^2\dd s<\infty\), so \(M\) is a true martingale. Consequently,
\[
\boxed{\E[X_t^2]=(x+\mu t)^2+\sigma^2t.}
\]
The answer is \((\text{mean})^2+\text{variance}\), as it should be.
\end{frame}

\begin{frame}{When may we erase the stochastic integral in expectation?}
After applying \(\Ito\), the random term is often
\[
\int_0^t H_s\dd W_s,
\qquad H_s=b_sf_x(s,X_s).
\]
If
\[
\E\!\left[\int_0^T H_s^2\dd s\right]<\infty,
\]
then this is an \(L^2\)-martingale and its expectation is zero.
\begin{alertblock}{The important distinction}
Pathwise local square-integrability generally gives only a \emph{local} martingale. Zero drift alone does not prove the martingale property.
\end{alertblock}
To exchange expectation and the drift integral, its absolute integrability must also be checked.
\end{frame}

\begin{frame}{A reusable log rule for positive processes}
Suppose \(X_t>0\) and
\[
\frac{\dd X_t}{X_t}=\alpha_t\dd t+\beta_t\dd W_t.
\]
With \(f(x)=\log x\),
\[
\boxed{
\dd\log X_t
=\left(\alpha_t-\frac12\beta_t^2\right)\dd t
+\beta_t\dd W_t.}
\]
\begin{block}{Volatility drag}
The arithmetic growth rate is \(\alpha_t\), but the log growth rate is reduced by \(\beta_t^2/2\).
\end{block}
The identity holds up to the first time \(X\) reaches zero if strict positivity is not known globally.
\end{frame}

\begin{frame}{Exercise: invert an Ito process}
Let \(X_t\neq0\) and
\[
\dd X_t=a_t\dd t+b_t\dd W_t.
\]
\begin{block}{Your turn}
\begin{enumerate}\setlength{\itemsep}{0.45em}
  \item Use \(f(x)=1/x\) to find \(\dd(1/X_t)\).
  \item Identify the term that the ordinary quotient rule would miss.
  \item Specialize to \(\dd X_t/X_t=\mu\dd t+\sigma\dd W_t\).
\end{enumerate}
\end{block}
Work only up to the first time \(X\) reaches zero.
\end{frame}

\begin{frame}{Exercise debrief: inversion creates a drift correction}
Since \(f'(x)=-x^{-2}\) and \(f''(x)=2x^{-3}\),
\[
\boxed{
\dd\!\left(\frac1{X_t}\right)
=\left(-\frac{a_t}{X_t^2}+\frac{b_t^2}{X_t^3}\right)\dd t
-\frac{b_t}{X_t^2}\dd W_t.}
\]
For geometric dynamics, \(R_t=1/X_t\) therefore satisfies
\[
\boxed{\frac{\dd R_t}{R_t}=(\sigma^2-\mu)\dd t-\sigma\dd W_t.}
\]
The extra \(+\sigma^2\) is the convexity correction. As a check, applying the product rule to \(X_tR_t=1\) makes every term cancel.
\end{frame}


% ================================================================
\subsection{Two-dimensional Ito formula and integration by parts}
% ================================================================

\begin{frame}{Two processes can share the same source of noise}
Consider two Ito processes driven by the same Brownian motion:
\[
\dd X_t=a_t^X\dd t+b_t^X\dd W_t,
\qquad
\dd Y_t=a_t^Y\dd t+b_t^Y\dd W_t.
\]
Their infinitesimal product satisfies
\[
\dd X_t\,\dd Y_t
=
b_t^Xb_t^Y(\dd W_t)^2
=
b_t^Xb_t^Y\dd t.
\]
Therefore,
\[
\boxed{\dd[X,Y]_t=b_t^Xb_t^Y\dd t.}
\]

\begin{block}{Interpretation}
If two processes are exposed to the same random shock, their fluctuations
create a non-zero cross-variation.
\end{block}

This cross term is what changes the usual multivariable chain rule.
\end{frame}


\begin{frame}{Exercise: what survives in a two-dimensional expansion?}
Suppose
\[
\dd X_t=a_t^X\dd t+b_t^X\dd W_t,
\qquad
\dd Y_t=a_t^Y\dd t+b_t^Y\dd W_t.
\]

\begin{block}{Your turn}
Using
\[
(\dd t)^2=0,
\qquad
\dd t\,\dd W_t=0,
\qquad
(\dd W_t)^2=\dd t,
\]
simplify:
\begin{enumerate}
    \item \((\dd X_t)^2\),
    \item \((\dd Y_t)^2\),
    \item \(\dd X_t\,\dd Y_t\).
\end{enumerate}
\end{block}

\begin{alertblock}{Question}
Why does the product \(\dd X_t\,\dd Y_t\) not disappear?
\end{alertblock}
\end{frame}


\begin{frame}{Ito's formula in two dimensions}
Let
\[
Z_t=f(t,X_t,Y_t),
\qquad f\in C^{1,2},
\]
where \(X\) and \(Y\) are continuous Ito processes. Then
\[
\boxed{
\begin{aligned}
\dd Z_t
={}&f_t\dd t
+f_x\dd X_t
+f_y\dd Y_t\\
&+\frac12 f_{xx}\dd[X]_t
+f_{xy}\dd[X,Y]_t
+\frac12 f_{yy}\dd[Y]_t.
\end{aligned}}
\]

\begin{block}{The new ingredient}
The mixed derivative \(f_{xy}\) is paired with the cross-variation
\(\dd[X,Y]_t\).
\end{block}

This is the stochastic version of the second-order Taylor expansion in two variables.
\end{frame}


\begin{frame}{Where does the mixed term come from?}
For a small change \((\dd X,\dd Y)\), the second-order Taylor expansion gives
\[
\begin{aligned}
\dd f
\approx{}&
f_x\dd X+f_y\dd Y\\
&+\frac12 f_{xx}(\dd X)^2
+f_{xy}\dd X\dd Y
+\frac12 f_{yy}(\dd Y)^2.
\end{aligned}
\]

In ordinary calculus, all quadratic terms vanish at first order.

In stochastic calculus,
\[
(\dd W)^2=\dd t,
\]
so quadratic terms involving Brownian increments survive.

\begin{block}{Key idea}
Ito's correction is simply the second-order Taylor expansion combined with
the special scaling \(\dd W\sim\sqrt{\dd t}\).
\end{block}
\end{frame}


\begin{frame}{Exercise: derive the stochastic product rule}
Take
\[
f(x,y)=xy.
\]
Then
\[
f_x=y,\qquad f_y=x,
\qquad f_{xx}=f_{yy}=0,
\qquad f_{xy}=1.
\]

\begin{block}{Your turn}
Apply the two-dimensional Ito formula to \(X_tY_t\).
What extra term appears compared with the ordinary product rule?
\end{block}

\vspace{0.3cm}

\begin{alertblock}{Hint}
Do not forget the mixed quadratic variation
\[
\dd[X,Y]_t.
\]
\end{alertblock}
\end{frame}


\begin{frame}{Integration by parts has one extra term}
Applying Ito's formula to \(f(x,y)=xy\) gives
\[
\boxed{
\dd(X_tY_t)
=
X_t\dd Y_t
+
Y_t\dd X_t
+
\dd[X,Y]_t.}
\]

Equivalently,
\[
\boxed{
X_tY_t
=
X_0Y_0
+\int_0^t X_s\dd Y_s
+\int_0^t Y_s\dd X_s
+[X,Y]_t.}
\]

\begin{block}{Compare with ordinary calculus}
The classical formula would contain only
\[
\dd(XY)=X\dd Y+Y\dd X.
\]
The quadratic covariation is the stochastic correction.
\end{block}
\end{frame}


\begin{frame}{Example: the product of two Ito processes}
Suppose
\[
\dd X_t=a_t^X\dd t+b_t^X\dd W_t,
\qquad
\dd Y_t=a_t^Y\dd t+b_t^Y\dd W_t.
\]
Since
\[
\dd[X,Y]_t=b_t^Xb_t^Y\dd t,
\]
integration by parts gives
\[
\begin{aligned}
\dd(X_tY_t)
={}&
\left(
X_ta_t^Y
+Y_ta_t^X
+b_t^Xb_t^Y
\right)\dd t\\
&+
\left(
X_tb_t^Y
+Y_tb_t^X
\right)\dd W_t.
\end{aligned}
\]

\begin{block}{What should you remember?}
Multiplying two stochastic processes creates an additional drift term:
\[
\boxed{b_t^Xb_t^Y\dd t.}
\]
\end{block}
\end{frame}


\begin{frame}{Independent Brownian noises have zero cross-variation}
Now suppose
\[
\dd X_t=a_t^X\dd t+b_t^X\dd W_t^1,
\qquad
\dd Y_t=a_t^Y\dd t+b_t^Y\dd W_t^2,
\]
where \(W^1\) and \(W^2\) are independent Brownian motions.

Then
\[
\dd W_t^1\dd W_t^2=0,
\]
and therefore
\[
\boxed{\dd[X,Y]_t=0.}
\]

\begin{block}{Consequently}
The product rule becomes
\[
\dd(X_tY_t)=X_t\dd Y_t+Y_t\dd X_t.
\]
There is no quadratic covariation correction.
\end{block}
\end{frame}


\begin{frame}{Correlated Brownian motions create cross terms}
Suppose instead
\[
\dd[W^1,W^2]_t=\rho\,\dd t,
\qquad -1\leq\rho\leq1.
\]
If
\[
\dd X_t=a_t^X\dd t+b_t^X\dd W_t^1,
\qquad
\dd Y_t=a_t^Y\dd t+b_t^Y\dd W_t^2,
\]
then
\[
\boxed{
\dd[X,Y]_t
=
\rho\,b_t^Xb_t^Y\dd t.}
\]

\begin{block}{Three useful cases}
\[
\rho=0:\ \text{no instantaneous covariance},
\qquad
\rho=1:\ \text{same shock},
\qquad
\rho=-1:\ \text{opposite shocks}.
\]
\end{block}
\end{frame}


\begin{frame}{Exercise: two correlated assets}
Suppose
\[
\frac{\dd S_t^1}{S_t^1}
=
\mu_1\dd t+\sigma_1\dd W_t^1,
\qquad
\frac{\dd S_t^2}{S_t^2}
=
\mu_2\dd t+\sigma_2\dd W_t^2,
\]
with
\[
\dd[W^1,W^2]_t=\rho\dd t.
\]

\begin{block}{Your turn}
\begin{enumerate}
    \item Compute \(\dd[S^1,S^2]_t\).
    \item Use integration by parts to find
    \(\dd(S_t^1S_t^2)\).
    \item What happens to the drift when \(\rho>0\)?
\end{enumerate}
\end{block}

\begin{alertblock}{Financial interpretation}
Correlation affects not only joint risk: through Ito's formula, it can also
appear as a drift correction after a nonlinear transformation.
\end{alertblock}
\end{frame}


% ================================================================
\subsection{Multidimensional Ito formula}
% ================================================================

\begin{frame}{A multidimensional diffusion has vectors and matrices}
Let
\[
X_t=(X_t^1,\ldots,X_t^d)^\top
\]
satisfy
\[
\boxed{
\dd X_t
=
b(t,X_t)\dd t
+
\sigma(t,X_t)\dd W_t,}
\]
where
\[
W_t\in\R^m,\qquad
b(t,x)\in\R^d,\qquad
\sigma(t,x)\in\R^{d\times m}.
\]

\begin{center}
\begin{tabular}{ccc}
\toprule
Object & Dimension & Role\\
\midrule
\(X_t\) & \(d\) & state variables\\
\(b\) & \(d\) & drift\\
\(W_t\) & \(m\) & independent shocks\\
\(\sigma\) & \(d\times m\) & exposure to shocks\\
\bottomrule
\end{tabular}
\end{center}
\end{frame}


\begin{frame}{The covariance matrix is \(\sigma\sigma^\top\)}
From
\[
\dd X_t=b_t\dd t+\sigma_t\dd W_t,
\]
the quadratic covariations satisfy
\[
\boxed{
\dd[X^i,X^j]_t
=
(\sigma_t\sigma_t^\top)_{ij}\dd t.}
\]

Define the instantaneous covariance matrix
\[
A_t=\sigma_t\sigma_t^\top.
\]

Then
\[
A_t^{ij}\dd t
\]
measures the instantaneous covariance between the \(i\)-th and \(j\)-th
components of \(X\).

\begin{block}{Important}
Even if the components of \(W\) are independent, the components of \(X\)
need not be independent: the matrix \(\sigma\) can mix the shocks.
\end{block}
\end{frame}


\begin{frame}{The multidimensional Ito formula}
Let
\[
\dd X_t=b_t\dd t+\sigma_t\dd W_t,
\qquad
f:[0,T]\times\R^d\to\R.
\]
For \(f\in C^{1,2}\),
\[
\boxed{
\begin{aligned}
\dd f(t,X_t)
={}&
\left(
\partial_t f
+
b_t\cdot\nabla f
+
\frac12
\sum_{i,j=1}^d
(\sigma_t\sigma_t^\top)_{ij}
\partial_{ij}f
\right)\dd t\\
&+
\nabla f(t,X_t)^\top\sigma_t\dd W_t.
\end{aligned}}
\]

\begin{block}{Same structure as in dimension one}
\[
\text{time}
+\text{drift}\times\text{slope}
+\frac12\,\text{covariance}\times\text{curvature}
+\text{noise}.
\]
\end{block}
\end{frame}


\begin{frame}{Matrix notation makes Ito's correction compact}
Let
\[
A_t=\sigma_t\sigma_t^\top
\]
and denote the Hessian of \(f\) by
\[
D^2f=
\begin{pmatrix}
\partial_{11}f & \cdots & \partial_{1d}f\\
\vdots & \ddots & \vdots\\
\partial_{d1}f & \cdots & \partial_{dd}f
\end{pmatrix}.
\]
Then
\[
\sum_{i,j=1}^d A_t^{ij}\partial_{ij}f
=
\operatorname{Tr}\!\left(A_tD^2f\right).
\]

Hence
\[
\boxed{
\dd f
=
\left(
f_t+b_t\cdot\nabla f
+\frac12\operatorname{Tr}(A_tD^2f)
\right)\dd t
+
\nabla f^\top\sigma_t\dd W_t.}
\]

\begin{block}{The object to remember}
In several dimensions,
\[
b^2f''\quad\longrightarrow\quad
\operatorname{Tr}(\sigma\sigma^\top D^2f).
\]
\end{block}
\end{frame}


\begin{frame}{Exercise: recover the two-dimensional formula}
Let
\[
X_t=
\begin{pmatrix}
X_t^1\\
X_t^2
\end{pmatrix},
\qquad
A_t=
\begin{pmatrix}
A_{11} & A_{12}\\
A_{12} & A_{22}
\end{pmatrix}.
\]

\begin{block}{Your turn}
Expand
\[
\frac12\operatorname{Tr}(A_tD^2f).
\]
Show that it equals
\[
\frac12 A_{11}f_{11}
+A_{12}f_{12}
+\frac12 A_{22}f_{22}.
\]
\end{block}

\begin{alertblock}{Why is there no \(1/2\) in front of the mixed term?}
The terms \((1,2)\) and \((2,1)\) both appear in the trace expansion,
and \(f_{12}=f_{21}\).
\end{alertblock}
\end{frame}


\begin{frame}{The generator is the drift produced by Ito's formula}
For the diffusion
\[
\dd X_t=b(t,X_t)\dd t+\sigma(t,X_t)\dd W_t,
\]
define
\[
\boxed{
\mathcal L f
=
b\cdot\nabla f
+
\frac12
\operatorname{Tr}
\left(\sigma\sigma^\top D^2f\right).}
\]
Then Ito's formula becomes
\[
\dd f(t,X_t)
=
\left(\partial_t f+\mathcal L f\right)(t,X_t)\dd t
+
\nabla f(t,X_t)^\top\sigma(t,X_t)\dd W_t.
\]

\begin{propositionblock}{A useful consequence}
If
\[
\partial_t f+\mathcal L f=0,
\]
then \(f(t,X_t)\) has zero drift and is a local martingale.
\end{propositionblock}
\end{frame}


% ================================================================
\section{Transformations Solve Canonical SDEs}
\subsection{GBM, mean reversion, and positive diffusions}
% ================================================================

\begin{frame}{Geometric Brownian motion models relative shocks}
Consider
\[
\dd S_t=\mu S_t\dd t+\sigma S_t\dd W_t,
\qquad S_0>0,\quad \sigma>0.
\]
Over a short interval,
\[
\frac{\Delta S_t}{S_t}\approx \mu\Delta t+\sigma\Delta W_t.
\]
Thus the same Brownian shock creates a larger dollar move when the asset level is larger.
\begin{block}{The useful transformation}
The coefficients are multiplicative, so take a logarithm to turn the SDE into an additive one.
\end{block}
\end{frame}

\begin{frame}{Ito solves geometric Brownian motion in one line}
The log rule gives
\[
\dd\log S_t
=\left(\mu-\frac12\sigma^2\right)\dd t+\sigma\dd W_t.
\]
Integrating from \(0\) to \(t\),
\[
\boxed{
S_t=S_0\exp\left\{\left(\mu-\frac12\sigma^2\right)t+\sigma W_t\right\}.}
\]
The explicit solution is strictly positive, so the logarithm was legitimate.
\begin{propositionblock}{Law}
\(\log S_t\) is Gaussian; therefore \(S_t\) is lognormal.
\end{propositionblock}
\end{frame}

\begin{frame}{Application: mean growth differs from typical growth}
From the explicit GBM solution,
\[
\E[S_t]=S_0e^{\mu t},
\qquad
\operatorname{Med}(S_t)=S_0e^{(\mu-\sigma^2/2)t}.
\]
For \(\sigma>0\),
\[
\Pp(S_T>K)
=\Pp\!\left(
Z>
\frac{\log(K/S_0)-(\mu-\sigma^2/2)T}{\sigma\sqrt T}
\right),
\quad Z\sim N(0,1).
\]
\begin{block}{Interpretation}
Volatility leaves the mean growth \(\mu\) unchanged but lowers the median log-growth. Rare large outcomes pull the mean upward.
\end{block}
\end{frame}

\begin{frame}{Worked finance example: one year of GBM in numbers}
Take
\[
S_0=100,\qquad \mu=8\%,\qquad \sigma=20\%,\qquad T=1.
\]
The explicit solution gives
\[
\E[S_1]=100e^{0.08}\approx108.33,
\qquad
\operatorname{Med}(S_1)=100e^{0.06}\approx106.18.
\]
For a target level of \(110\),
\[
\Pp(S_1>110)
=1-\Phi\!\left(\frac{\log(1.10)-0.06}{0.20}\right)
\approx\boxed{43.0\%}.
\]
\begin{block}{Reading the numbers}
The mean, median, and probability of beating a target answer three different questions. A positive expected return does not remove downside uncertainty.
\end{block}
\end{frame}

\begin{frame}{Exercise: powers of a geometric Brownian motion}
Let \(p\in\R\) and \(Y_t=S_t^p\).
\begin{block}{Your turn}
\begin{enumerate}
\item Find \(\dd Y_t/Y_t\).
\item Choose \(\kappa\) so that \(e^{-\kappa t}S_t^p\) has zero drift.
\item Check your formula for \(p=0,1,2\).
\end{enumerate}
\end{block}
The cases \(p=-1\) and \(p=2\) are particularly useful consistency checks.
\end{frame}

\begin{frame}{Exercise debrief: curvature changes the growth rate}
Since \(f'(s)=ps^{p-1}\) and \(f''(s)=p(p-1)s^{p-2}\),
\[
\boxed{
\frac{\dd(S_t^p)}{S_t^p}
=\left[p\mu+\frac12p(p-1)\sigma^2\right]\dd t
+p\sigma\dd W_t.}
\]
Hence
\[
\kappa=p\mu+\frac12p(p-1)\sigma^2.
\]
For constant coefficients, \(e^{-\kappa t}S_t^p\) is lognormal with constant mean and is a true martingale.
\begin{block}{Why the correction changes sign}
Powers are convex for \(p>1\) and concave for \(0<p<1\); the \(\Ito\) correction detects this curvature.
\end{block}
\end{frame}

\begin{frame}{Time-dependent GBM uses the same transformation}
Suppose
\[
\frac{\dd S_t}{S_t}=\mu_t\dd t+\sigma_t\dd W_t,
\]
with progressively measurable coefficients satisfying
\[
\int_0^T|\mu_t|\dd t<\infty,
\qquad
\int_0^T\sigma_t^2\dd t<\infty
\quad\text{a.s.}
\]
Then
\[
\boxed{
S_t=S_0\exp\left\{
\int_0^t\left(\mu_s-\frac12\sigma_s^2\right)\dd s
+\int_0^t\sigma_s\dd W_s
\right\}.}
\]
If \(\mu\) and \(\sigma\) are deterministic, the conditional law remains lognormal.
\end{frame}

\begin{frame}{Mean reversion: the Ornstein--Uhlenbeck model}
The SDE
\[
\dd X_t=\kappa(m-X_t)\dd t+\sigma\dd W_t,
\qquad \kappa>0,
\]
has a restoring drift:
\begin{itemize}
\item above \(m\), the drift is negative;
\item below \(m\), the drift is positive;
\item \(1/\kappa\) is the time scale of mean reversion.
\end{itemize}
It appears in short-rate models, spreads, temperatures, and other persistent but mean-reverting quantities.
\end{frame}

\begin{frame}{A deterministic integrating factor removes the drift}
Rewrite the SDE as
\[
\dd X_t+\kappa X_t\dd t=\kappa m\dd t+\sigma\dd W_t.
\]
Because \(e^{\kappa t}\) has finite variation,
\[
\dd(e^{\kappa t}X_t)
=e^{\kappa t}\dd X_t+\kappa e^{\kappa t}X_t\dd t,
\]
with no covariation correction. Therefore
\[
\dd(e^{\kappa t}X_t)
=\kappa m e^{\kappa t}\dd t+\sigma e^{\kappa t}\dd W_t.
\]
Integrating and dividing by \(e^{\kappa t}\) solves the SDE.
\end{frame}

\begin{frame}{The OU solution displays memory loss explicitly}
For \(0\le s<t\),
\[
\boxed{
X_t=m+(X_s-m)e^{-\kappa(t-s)}
+\sigma\int_s^t e^{-\kappa(t-u)}\dd W_u.}
\]
Conditionally on \(\F_s\),
\[
\E[X_t\mid\F_s]=m+(X_s-m)e^{-\kappa(t-s)},
\]
and
\[
\Var(X_t\mid\F_s)
=\frac{\sigma^2}{2\kappa}\left(1-e^{-2\kappa(t-s)}\right).
\]
The effect of the current deviation \(X_s-m\) decays exponentially.
\end{frame}

\begin{frame}{Exercise: read the long-run behavior from the formula}
For the OU process, determine:
\begin{block}{Your turn}
\begin{enumerate}\setlength{\itemsep}{0.55em}
  \item the limiting conditional mean;
  \item the limiting conditional variance;
  \item the half-life of a deviation from \(m\).
\end{enumerate}
\end{block}
Then evaluate the half-life when \(\kappa=0.7\) per year.
\end{frame}

\begin{frame}{Exercise debrief: mean reversion has a measurable speed}
As \(t-s\to\infty\),
\[
\E[X_t\mid\F_s]\to m,
\qquad
\Var(X_t\mid\F_s)\to\frac{\sigma^2}{2\kappa}.
\]
The equation \(e^{-\kappa h}=1/2\) gives the half-life
\[
\boxed{h=\frac{\log2}{\kappa}.}
\]
For \(\kappa=0.7\),
\[
h\approx0.99\text{ year},
\qquad
e^{-0.7\times2}\approx24.7\%.
\]
After two years, only about one quarter of today's deviation from \(m\) remains in the conditional mean.
\end{frame}

\begin{frame}{Exercise: obtain the CIR mean without solving the SDE}
Assume a nonnegative CIR process satisfies
\[
\dd R_t=(\alpha-\beta R_t)\dd t+\sigma\sqrt{R_t}\dd W_t,
\qquad R_0=r_0,\quad \alpha,\beta,\sigma>0.
\]
\begin{block}{Your turn}
\begin{enumerate}\setlength{\itemsep}{0.5em}
  \item Compute \(\dd(e^{\beta t}R_t)\).
  \item Integrate and take expectations.
  \item Find \(\E[R_t]\) and its long-run limit.
\end{enumerate}
\end{block}
What part of the calculation uses the square-root volatility, and what part does not?
\end{frame}

\begin{frame}{Exercise debrief: CIR and Vasicek share the same mean}
With \(X_t=e^{\beta t}R_t\), the product rule gives
\[
\dd X_t
=\alpha e^{\beta t}\dd t
+\sigma e^{\beta t}\sqrt{R_t}\dd W_t.
\]
After localization and the usual integrability check, the stochastic integral has mean zero. Therefore
\[
\E[X_t]=r_0+\frac{\alpha}{\beta}(e^{\beta t}-1),
\]
and
\[
\boxed{
\E[R_t]
=r_0e^{-\beta t}+\frac{\alpha}{\beta}(1-e^{-\beta t})
\longrightarrow\frac{\alpha}{\beta}.}
\]
\begin{block}{Interpretation}
The diffusion coefficient controls dispersion, but the first moment follows the same mean-reversion ODE as the Vasicek model. The second moment is a natural next exercise.
\end{block}
\end{frame}

% ================================================================
\section{Products, Ratios, and Several Risk Factors}
\subsection{Mixed covariation, products, and hedging}
% ================================================================

\begin{frame}{The two-dimensional Ito formula includes a mixed term}
For \(f\in C^{1,2}([0,T]\times\R^2)\) and continuous Ito processes \(X,Y\),
\[
\begin{aligned}
\dd f(t,X_t,Y_t)
={}&f_t\dd t+f_x\dd X_t+f_y\dd Y_t\\
&+\frac12f_{xx}\dd[X]_t
+f_{xy}\dd[X,Y]_t
+\frac12f_{yy}\dd[Y]_t.
\end{aligned}
\]
If the driving Brownian motions satisfy \(\dd W_t\dd B_t=\rho\dd t\), the mixed term carries the correlation.
\begin{block}{Memory aid}
Diagonal second derivatives receive a factor \(1/2\); the mixed derivative does not, because it appears twice in the Hessian sum.
\end{block}
\end{frame}

\begin{frame}{The stochastic product rule has one extra term}
Choosing \(f(x,y)=xy\) gives
\[
\boxed{\dd(X_tY_t)=X_t\dd Y_t+Y_t\dd X_t+\dd[X,Y]_t.}
\]
For
\[
\dd X_t=a_t\dd t+b_t\dd W_t,
\qquad
\dd Y_t=c_t\dd t+d_t\dd B_t,
\]
with \(\dd W_t\dd B_t=\rho\dd t\),
\[
\dd[X,Y]_t=\rho b_td_t\dd t.
\]
This correction is the continuous-time record of simultaneous random shocks.
\end{frame}

\begin{frame}{Application: the product of two correlated assets}
Suppose
\[
\frac{\dd S_t}{S_t}=\mu_S\dd t+\sigma_S\dd W_t,
\qquad
\frac{\dd V_t}{V_t}=\mu_V\dd t+\sigma_V\dd B_t,
\]
and \(\dd W_t\dd B_t=\rho\dd t\). Then
\[
\boxed{
\frac{\dd(S_tV_t)}{S_tV_t}
=\left(\mu_S+\mu_V+\rho\sigma_S\sigma_V\right)\dd t
+\sigma_S\dd W_t+\sigma_V\dd B_t.}
\]
Positive correlation increases the product's drift through covariation; negative correlation decreases it.
\end{frame}

\begin{frame}{Exercise: normalize the sum of two Brownian motions}
Let \(W,B\) be jointly Brownian with constant correlation \(\rho>-1\).
\begin{block}{Your turn}
\begin{enumerate}
\item Compute \(\dd[(W+B)]_t\).
\item Find a constant \(c\) such that \(c(W+B)\) has quadratic variation \(t\).
\item What happens when \(\rho=-1\)?
\end{enumerate}
\end{block}
\end{frame}

\begin{frame}{Exercise debrief: correlation changes total variance}
Using the product table,
\[
\dd[W+B]_t
=(\dd W_t+\dd B_t)^2
=2(1+\rho)\dd t.
\]
Therefore
\[
\boxed{\widetilde W_t=\frac{W_t+B_t}{\sqrt{2(1+\rho)}}}
\]
is a standard Brownian motion. This follows from joint Gaussianity, or from Levy's characterization.
If \(\rho=-1\), then \(B=-W\) up to indistinguishability and the sum is identically zero.
\end{frame}

\begin{frame}{Exercise: derive the ratio of two risky assets}
With the same dynamics for \(S\) and \(V\), let \(R_t=S_t/V_t\).
\begin{block}{Your turn}
Use either the two-dimensional \(\Ito\) formula or the product rule applied to \(S_t(1/V_t)\). Find \(\dd R_t/R_t\).
\end{block}
\begin{alertblock}{Common trap}
The ordinary quotient rule misses both the inverse correction \(\sigma_V^2\dd t\) and the cross correction \(-\rho\sigma_S\sigma_V\dd t\).
\end{alertblock}
\end{frame}

\begin{frame}{Exercise debrief: the denominator contributes convexity}
The ratio satisfies
\[
\boxed{
\frac{\dd R_t}{R_t}
=\left(
\mu_S-\mu_V+\sigma_V^2-\rho\sigma_S\sigma_V
\right)\dd t
+\sigma_S\dd W_t-\sigma_V\dd B_t.}
\]
Sanity checks:
\begin{itemize}
\item if \(S=V\), then \(\rho=1\), \(\mu_S=\mu_V\), \(\sigma_S=\sigma_V\), and every term vanishes;
\item if \(V\) is deterministic, \(\sigma_V=0\), and ordinary quotient calculus returns.
\end{itemize}
\end{frame}

\begin{frame}{Application: correlation determines a minimum-variance hedge}
Consider the instantaneous shock
\[
\sigma_X\dd W_t-h\sigma_Y\dd B_t,
\qquad \dd W_t\dd B_t=\rho\dd t.
\]
Its variance rate is
\[
q(h)=\sigma_X^2+h^2\sigma_Y^2-2h\rho\sigma_X\sigma_Y.
\]
Minimizing this quadratic gives
\[
\boxed{h^*=\rho\frac{\sigma_X}{\sigma_Y},
\qquad q(h^*)=\sigma_X^2(1-\rho^2).}
\]
Only the correlated component can be hedged; orthogonal noise remains.
\begin{block}{Numerical check}
If \(\sigma_X=25\%\), \(\sigma_Y=20\%\), and \(\rho=0.8\), then
\[
h^*=1,
\qquad
\sqrt{q(h^*)}=25\%\sqrt{1-0.8^2}=15\%.
\]
\end{block}
\end{frame}

\begin{frame}{Optional extension: rotating shocks preserves Brownian risk}
Let \(W\) be a standard \(d\)-dimensional Brownian motion and let \(H_t\) be predictable with
\[
H_tH_t^\top=I_d.
\]
Define
\[
B_t=\int_0^t H_s\dd W_s.
\]
Then \(B\) is a continuous local martingale and
\[
\dd[B^i,B^j]_t=(H_tH_t^\top)_{ij}\dd t=\delta_{ij}\dd t.
\]
By Levy's characterization, \(B\) is again a standard \(d\)-dimensional Brownian motion.
\begin{block}{One-dimensional surprise}
\(\int_0^t\operatorname{sgn}(W_s)\dd W_s\) is Brownian, although its integrand depends on the path.
\end{block}
\end{frame}

% ================================================================
\section{Finance Applications of Ito Calculus}
\subsection{Discounting, delta hedging, and realized variance}
% ================================================================

\begin{frame}{Exercise: what does discounting do under the physical measure?}
Let the bank account and asset satisfy
\[
\dd A_t=rA_t\dd t,
\qquad
\dd S_t=\mu S_t\dd t+\sigma S_t\dd W_t^{\Pp}.
\]
Set \(\widetilde S_t=S_t/A_t\).
\begin{block}{Your turn}
Find \(\dd\widetilde S_t/\widetilde S_t\). Does discounting by the bank account automatically make \(\widetilde S\) a martingale under \(\Pp\)?
\end{block}
Hint: \(A_t^{-1}=A_0^{-1}e^{-rt}\) has finite variation.
\end{frame}

\begin{frame}{Exercise debrief: the measure matters}
The product rule gives
\[
\boxed{
\frac{\dd\widetilde S_t}{\widetilde S_t}
=(\mu-r)\dd t+\sigma\dd W_t^{\Pp}.}
\]
Under the physical measure \(\Pp\), the drift generally remains \(\mu-r\).
Under a pricing measure \(\mathbb Q\), the model is written
\[
\frac{\dd S_t}{S_t}=r\dd t+\sigma\dd W_t^{\mathbb Q},
\]
so \(\widetilde S\) has zero drift and is a local martingale. Standard integrability conditions make it a true martingale.
\end{frame}

\begin{frame}{Ito turns an option value into drift plus hedgeable risk}
Suppose \(V_t=v(t,S_t)\), with
\[
\dd S_t=\mu S_t\dd t+\sigma S_t\dd W_t.
\]
Then
\[
\boxed{
\dd V_t
=\left(v_t+\mu S_tv_s+\frac12\sigma^2S_t^2v_{ss}\right)\dd t
+\sigma S_tv_s\dd W_t.}
\]
The derivatives have financial names:
\[
\Delta=v_s,
\qquad
\Gamma=v_{ss}.
\]
Delta controls first-order exposure; gamma measures curvature and creates the \(\Ito\) correction.
\end{frame}

\begin{frame}{Delta hedging cancels the instantaneous Brownian shock}
\small
Let \(\Delta_t=v_s(t,S_t)\), and let the cash account finance every change in the stock position. The self-financing hedged \emph{gain} is
\[
\dd\Pi_t:=\dd V_t-\Delta_t\dd S_t.
\]
Substitution gives
\[
\begin{aligned}
\dd\Pi_t
&=\left(v_t+\mu S_tv_s+\frac12\sigma^2S_t^2v_{ss}\right)\dd t
+\sigma S_tv_s\dd W_t\\
&\quad-v_s\left(\mu S_t\dd t+\sigma S_t\dd W_t\right)\\
&=\boxed{\left(v_t+\frac12\sigma^2S_t^2v_{ss}\right)\dd t.}
\end{aligned}
\]
The unknown physical drift \(\mu\) cancels together with the random shock.
\begin{alertblock}{Bookkeeping point}
This is the gain of a dynamically financed strategy. It is not the algebraic identity
\(\dd(V_t-\Delta_tS_t)=\dd V_t-\Delta_t\dd S_t\).
\end{alertblock}
\end{frame}

\begin{frame}{Application: curvature creates deterministic hedged P\&L}
\small
Take the toy claim \(V_t=S_t^2\). Then
\[
\Delta_t=2S_t,
\qquad
\Gamma_t=2.
\]
The \(\Ito\) formula gives
\[
\dd(S_t^2)=2S_t\dd S_t+\sigma^2S_t^2\dd t.
\]
After shorting \(2S_t\) shares,
\[
\boxed{\dd V_t-2S_t\dd S_t=\sigma^2S_t^2\dd t.}
\]
Continuous rebalancing converts quadratic price variation into a finite-variation term.
\begin{block}{One-day numerical scale}
For an option with \(\Gamma=0.02\), \(S=100\), \(\sigma=20\%\), and \(\Delta t=1/252\), the gamma contribution is
\[
\frac12\Gamma\sigma^2S^2\Delta t\approx\$0.016
\]
per option, before theta and financing.
\end{block}
\end{frame}

\begin{frame}{Preview for Session 5: the Black--Scholes PDE}
If the delta-hedged position is self-financing and locally riskless, it must earn the bank-account rate on its value \(V_t-\Delta_tS_t\):
\[
\dd V_t-\Delta_t\dd S_t
=r(V_t-\Delta_tS_t)\dd t.
\]
Substituting \(\Delta=v_s\) gives
\[
\boxed{
v_t+\frac12\sigma^2s^2v_{ss}+rsv_s-rv=0.}
\]
\begin{block}{Preview}
Quadratic variation produces the gamma term; no-arbitrage determines how that locally riskless term must grow.
\end{block}
\end{frame}

\begin{frame}{Exercise: convexity, concavity, and the Ito correction}
Let \(\dd X_t=a_t\dd t+b_t\dd W_t\) and let \(f\) have no explicit time dependence.
\begin{block}{Your turn}
Compare the drift of \(f(X_t)\) with the first-order prediction \(f'(X_t)a_t\dd t\) when:
\begin{enumerate}
\item \(f\) is convex;
\item \(f\) is concave;
\item \(f\) is linear.
\end{enumerate}
Use \(f(x)=x^2\) and \(f(x)=\log x\) as examples.
\end{block}
\end{frame}

\begin{frame}{Exercise debrief: volatility interacts with curvature}
The correction is
\[
\frac12b_t^2f''(X_t)\dd t.
\]
Therefore:
\begin{center}
\begin{tabular}{lcc}
\toprule
Shape of \(f\) & Sign of \(f''\) & Volatility correction \\
\midrule
Convex & \(f''\ge0\) & nonnegative \\
Linear & \(f''=0\) & zero \\
Concave & \(f''\le0\) & nonpositive \\
\bottomrule
\end{tabular}
\end{center}
Thus volatility raises the drift of a convex transform and lowers that of a concave transform. This local fact is a cornerstone of option convexity.
\end{frame}

\begin{frame}{Application: log quadratic variation is realized variance}
\small
For a positive asset with
\[
\frac{\dd S_t}{S_t}=\mu_t\dd t+\sigma_t\dd W_t,
\qquad
\dd\log S_t
=\left(\mu_t-\frac12\sigma_t^2\right)\dd t+\sigma_t\dd W_t.
\]
Therefore
\[
\boxed{[\log S]_T=\int_0^T\sigma_t^2\dd t.}
\]
On a fine grid, squared log returns estimate this integrated variance.
\begin{block}{Four-day calculation}
For log returns \(1.0\%,-0.5\%,0.7\%,-0.2\%\),
\[
\widehat{RV}=0.01^2+(-0.005)^2+0.007^2+(-0.002)^2
=1.78\times10^{-4},
\qquad
\widehat\sigma_{\rm ann}=\sqrt{\frac{252}{4}\widehat{RV}}\approx10.6\%.
\]
\end{block}
\end{frame}

% ================================================================
\section{Guided Practice and Common Traps}
\subsection{Exercises, diagnostics, and checklist}
% ================================================================

\begin{frame}{Rapid practice I: three transformations}
Compute the differential of each process.
\begin{enumerate}
\item \(Y_t=e^{\alpha t+\beta W_t}\).
\item \(Z_t=W_tB_t\), where \(\dd W_t\dd B_t=\rho\dd t\).
\item Use \(\Ito\) on \(e^{W_t}\) to remove \(\dd W\) from
\(\int_0^t e^{W_s}\dd W_s\).
\end{enumerate}
\begin{block}{Work pattern}
Write the relevant derivatives, substitute the dynamics, use the product table, and finally group drift and diffusion.
\end{block}
\end{frame}

\begin{frame}{Rapid practice I: answers}
\begin{align*}
\dd Y_t
&=Y_t\left(\alpha+\frac12\beta^2\right)\dd t
+\beta Y_t\dd W_t,\\
\dd Z_t
&=W_t\dd B_t+B_t\dd W_t+\rho\dd t,\\
\int_0^t e^{W_s}\dd W_s
&=e^{W_t}-1-\frac12\int_0^t e^{W_s}\dd s.
\end{align*}
Checks:
\begin{itemize}
\item the exponential is convex, so its correction is positive;
\item correlation contributes finite variation to the product;
\item integrating \(\dd(e^{W_t})\) isolates the stochastic integral.
\end{itemize}
\end{frame}

\begin{frame}{Rapid practice II: manufacture a martingale}
Let
\[
M_t=e^{at}\cos(bW_t),
\qquad a,b\in\R.
\]
\begin{block}{Your turn}
Choose \(a\) so that \(M\) has zero drift. Then decide whether it is only a local martingale or a true martingale on \([0,T]\).
\end{block}
\vspace{0.3em}
Do not stop after saying ``the drift is zero.'' Check the stochastic integrand.
\end{frame}

\begin{frame}{Rapid practice II: answer and integrability check}
The \(\Ito\) formula gives
\[
\dd M_t
=e^{at}\left(a-\frac12b^2\right)\cos(bW_t)\dd t
-be^{at}\sin(bW_t)\dd W_t.
\]
Thus \(a=b^2/2\) kills the drift. On \([0,T]\),
\[
\E\int_0^T b^2e^{b^2t}\sin^2(bW_t)\dd t
\le b^2\int_0^T e^{b^2t}\dd t<\infty.
\]
Therefore the stochastic integral is in \(L^2\), and \(M\) is a true martingale on every bounded horizon.
\end{frame}

\begin{frame}{Error diagnosis: find every incorrect line}
A student writes, for \(\dd X_t=a_t\dd t+b_t\dd W_t\),
\[
\begin{aligned}
\dd f(t,X_t)
&=f_x\dd X_t+\frac12b_tf_{xx}\dd t,\\
\dd(W_tB_t)&=W_t\dd B_t+B_t\dd W_t,\\
\text{zero drift}&\Longrightarrow\text{true martingale},\\
\frac{\dd(S_t/A_t)}{S_t/A_t}
&=\sigma\dd W_t^{\Pp}.
\end{aligned}
\]
Here \(\dd W_t\dd B_t=\rho\dd t\), \(\dd S_t/S_t=\mu\dd t+\sigma\dd W_t^{\Pp}\), and \(\dd A_t/A_t=r\dd t\).
\begin{block}{Your turn}
Correct the four lines and name the missing hypothesis or term in each one.
\end{block}
\end{frame}

\begin{frame}{Error diagnosis: corrections}
The corrected statements are
\[
\begin{aligned}
\dd f(t,X_t)
&=\left(f_t+a_tf_x+\frac12b_t^2f_{xx}\right)\dd t+b_tf_x\dd W_t,\\
\dd(W_tB_t)&=W_t\dd B_t+B_t\dd W_t+\rho\dd t,\\
\text{zero drift}&\Longrightarrow\text{local martingale first},\\
\frac{\dd(S_t/A_t)}{S_t/A_t}
&=(\mu-r)\dd t+\sigma\dd W_t^{\Pp}.
\end{aligned}
\]
To obtain a true martingale, verify integrability. Discounting kills the drift only after moving to a pricing measure under which the asset drift is \(r\).
\end{frame}

\begin{frame}{Optional preview: when the classical Ito formula is not enough}
The formula used today assumes \(f\in C^{1,2}\). It cannot be applied directly at the kink of
\[
f(x)=|x|
\qquad\text{or}\qquad
f(s)=(s-K)^+.
\]
For nonsmooth convex functions, a local-time term appears in the generalized formula.
\begin{alertblock}{Practical rule}
Before differentiating, locate singularities and ask whether the process can reach them. If it can, smooth approximation or the Ito--Tanaka formula may be needed.
\end{alertblock}
This matters for option payoffs even when the option value itself becomes smooth before maturity.
\end{frame}

\begin{frame}{A compact checklist for every Ito calculation}
\begin{enumerate}
\item Identify time and state variables; compute \(f_t,f_x,f_{xx}\).
\item Write the driving dynamics and their covariations.
\item Apply the product table: \((\dd W)^2=\dd t\), \(\dd W\dd B=\rho\dd t\).
\item Keep every term of order \(\dd t\); discard higher-order products.
\item Group the result into drift and diffusion.
\item Check the domain and regularity of \(f\).
\item If claiming a martingale or taking expectations, verify integrability.
\item In finance, state the probability measure.
\end{enumerate}
\end{frame}

\begin{frame}{What to remember from Session 4}
\begin{itemize}
\item Brownian quadratic variation makes stochastic calculus second order.
\item Covariation turns correlation into a \(\dd t\)-term.
\item The \(\Ito\) formula is a chain rule with a curvature correction.
\item Transformations solve useful SDEs and moment equations for GBM, OU, and CIR.
\item Products, ratios, discounting, and delta hedging all use the same small table.
\item Zero drift means local martingale; suitable integrability upgrades the conclusion.
\end{itemize}
\begin{propositionblock}{Core habit}
Always separate drift, diffusion, covariation, and measure. Most mistakes disappear once those four ingredients are visible.
\end{propositionblock}
\end{frame}

\end{document}
